% Optional IEEE conference short paper (4--8 pages stub)
% Compile: pdflatex ieee_short_paper && bibtex ieee_short_paper && pdflatex x2
% Overleaf: upload this file + references.bib in same project
\documentclass[conference]{IEEEtran}
\IEEEoverridecommandlockouts

\usepackage[T1]{fontenc}
\usepackage[utf8]{inputenc}
\usepackage{cite}
\usepackage{amsmath,amssymb}
\usepackage{graphicx}
\usepackage{booktabs}
\usepackage{url}
\usepackage{hyperref}

\hypersetup{hidelinks}

\begin{document}

\title{Augmented Analytics on a Production Lakehouse:\\Integrating Generative AI into European Job Intelligence}

\author{\IEEEauthorblockN{Ritesh Rakesh Jadhav}
\IEEEauthorblockA{\textit{M.Sc.\ Data Science}\\
\textit{University of Europe for Applied Sciences}\\
Potsdam, Germany\\
\texttt{[student.email@ue-germany.de]}}}

\maketitle

\begin{abstract}
European labour-market data is fragmented across government portals, aggregators, and company applicant-tracking systems.
DataForge is a live serverless medallion lakehouse on AWS that ingests multi-source vacancies, preserves slowly changing history, and serves public analytics APIs.
This paper describes a design-science integration of generative AI---model gateway, rules-first enrichment, hybrid retrieval Match API, and reproducible evaluation harness---without mutating dimensional keys.
On a labelled set of 103 queries over 96 jobs, dense embedding retrieval achieves nDCG@10 of 0.217 compared with 0.146 for a rule-based Career Matching Wizard and 0.135 for BM25 alone.
A modelled rules-first enrichment path costs approximately \EUR{0.06} per thousand jobs.
Live multi-provider Pareto experiments on AWS Bedrock remain pending account quota approval; architecture, baselines, and partial results are reported with explicit operational limits.
\end{abstract}

\begin{IEEEkeywords}
lakehouse, generative AI, information retrieval, design science, EU labour market, LLMOps
\end{IEEEkeywords}

\section{Introduction}
\label{sec:intro}

Job intelligence platforms must combine reliable data engineering with emerging AI capabilities.
In Europe, vacancy feeds originate from the German Federal Employment Agency (BA Jobsuche), EURES, aggregators such as Arbeitnow, and numerous company ATS endpoints.
Production systems require ingest validation, historical versioning, cost control, and compliance-aware serving layers---requirements that notebook-scale retrieval demos often omit.

DataForge addresses the engineering substrate: daily Bronze snapshots, Silver SCD Type 2 history, Gold analytics, API Gateway endpoints, and a public dashboard.
Industry demand has shifted from raw CSV exports toward \emph{augmented analytics}---semantic matching, structured enrichment, and measurable unit economics.
This work asks how generative AI can be added to such a pipeline under EU residency preferences, free-tier budgets, and falsifiable baselines.

Contributions:
\begin{itemize}
  \item a reference architecture for additive GenAI on medallion layers;
  \item a multi-provider gateway with kill switch, schema validation, and cost logging;
  \item an labelled IR evaluation comparing dense, hybrid, BM25, and heuristic matchers;
  \item a modelled unit-cost analysis for rules-first versus full-LLM enrichment.
\end{itemize}

\section{Related Work and Gap}
\label{sec:related}

Medallion lakehouses unify warehousing and ML-oriented storage~\cite{armbrust2021lakehouse}; Kimball-style SCD Type 2 preserves job posting history~\cite{kimball2013dwt}.
Dense passage retrieval~\cite{karpukhin2020dpr} and retrieval-augmented generation~\cite{lewis2020rag} motivate embedding-based ranking over keyword heuristics.
Graded ranking quality is measured with nDCG~\cite{jarvelin2002ndcg}.
Design science research frames artefact-building theses that must still ground claims in theory~\cite{hevner2004dsr}.

Few applied studies combine (i) a live multi-source lakehouse, (ii) controlled multi-provider GenAI integration, (iii) an honest non-AI baseline, and (iv) unit-cost grounding.
DataForge provides that combined setting.

\section{System and Method}
\label{sec:method}

\subsection{Architecture}

Public job documents flow Bronze $\rightarrow$ Silver (SCD2) $\rightarrow$ Gold and APIs.
Generative components are \emph{consumers}: rules-first enrichment and hybrid Match API read Silver/Gold without rewriting \texttt{job\_id} or validity intervals.
A \texttt{ModelRouter} routes tasks (\texttt{enrich}, \texttt{embed}, \texttt{explain}) across Bedrock (EU), OpenAI, Anthropic, and local providers with JSON schema checks, prompt versioning, daily budget caps, and an \texttt{AI\_ENABLED} kill switch.

\subsection{Matching pipeline}

Hybrid retrieval fuses BM25 lexical scores with dense embeddings via reciprocal rank fusion.
The Career Matching Wizard baseline applies fixed keyword weights (60\% skills, 20\% seniority, 20\% location).

\subsection{Evaluation design}

RQ3 uses 42 author-graded queries against 96 jobs; primary metric nDCG@10, with P@5, Recall@20, and MRR as secondary measures.
RQ2 compares provider quality--cost--latency Pareto points; rules-only baseline is complete; live Bedrock Nova Micro runs were blocked at zero account quota at measurement time.
RQ4 models enrichment, embedding, and match cost for a 1{,}000-job scenario using public list prices and a 30\% LLM-call share under rules-first enrichment.

\section{Results}
\label{sec:results}

Table~\ref{tab:rq3} summarises matching quality.
Dense retrieval leads on nDCG@10 (0.217), supporting hypothesis H1 on this label set.
Hybrid retrieval (0.171) remains the product default for robustness across lexical and semantic signals despite lower nDCG on this gold set.

\begin{table}[ht]
\centering
\caption{Matching quality on labelled set (103 queries, 96 jobs). Values may update as labels grow.}
\label{tab:rq3}
\begin{tabular}{@{}lrrrr@{}}
\toprule
Method & nDCG@10 & P@5 & R@20 & MRR \\
\midrule
BM25 & 0.135 & 0.162 & 0.214 & 0.184 \\
Dense & \textbf{0.217} & \textbf{0.219} & \textbf{0.274} & \textbf{0.268} \\
Hybrid & 0.171 & 0.171 & 0.234 & 0.234 \\
Heuristic wizard & 0.146 & 0.191 & 0.240 & 0.212 \\
\bottomrule
\end{tabular}
\end{table}

Integration checks (RQ1) confirm additive schemas, operational kill switch, and unchanged SCD semantics; nightly Bedrock enrichment was paused pending quota increase.
Unit-cost modelling (RQ4) yields approximately \EUR{0.06} per thousand jobs on the rules-first path (enrichment sample 30\%, embeddings, 100 match requests).

\section{Discussion and Limitations}
\label{sec:discussion}

Results support three practical implications: enforce rules-before-tokens enrichment, treat AI outputs as additive artefacts, and always benchmark against simple baselines.
Compliance framing treats Match as decision support for job seekers on public vacancy text, not automated hiring; GDPR and EU AI Act documentation remains necessary.

Limitations include author labels, a single cloud region, modest absolute nDCG, and incomplete live RQ2 provider cells.
Absolute scores should not be compared to open-domain QA benchmarks.

\section{Conclusion}
\label{sec:conclusion}

This case study shows that generative AI can extend a production European job lakehouse when integration is constrained by historical keys, evaluation harnesses, and cost controls.
Dense retrieval outperforming a transparent heuristic wizard demonstrates measurable value for augmented analytics claims.
Future work includes completing Bedrock Pareto experiments after quota approval, enlarging the gold label set with peer review, and flushing live cost logs from production Lambdas.

\section*{Acknowledgment}

The author thanks proposed supervisors at UE Potsdam and the open-source DataForge contributors.

\bibliographystyle{IEEEtran}
\bibliography{references}

\end{document}
